\section{Fundamentação Teórica}
\label{sec:theory}

\subsection{Grafos, emparelhamentos e cobertura por vértices}

Um grafo não direcionado $G=(V,E)$ possui um conjunto de vértices $V$ e um
conjunto de arestas $E$, em que cada aresta é um par não ordenado de
vértices. Neste trabalho, os grafos são simples: não há laços nem arestas
paralelas. O grau $d_G(v)$ é a quantidade de arestas incidentes em $v$;
um vértice de grau zero é isolado. Um grafo completo $K_n$ conecta todos
os pares distintos. Um grafo bipartido admite uma partição
$V=A\mathbin{\dot\cup}B$ sem arestas internas às partes; no bipartido
completo $K_{a,b}$ existem todas as $ab$ arestas entre elas
\cite{cormen2022}.

Um emparelhamento (\emph{matching}) $M\subseteq E$ é um conjunto de arestas
que não compartilham extremos. Ele é \emph{maximal} quando nenhuma aresta
adicional pode ser incluída mantendo essa propriedade. Isso difere de
um emparelhamento \emph{máximo}, que possui a maior cardinalidade entre
todos os emparelhamentos do grafo. A maximalidade é suficiente para a
aproximação implementada, embora não implique cardinalidade máxima
\cite{vazirani2001}.

Uma cobertura por vértices é um conjunto $C\subseteq V$ tal que
\begin{equation}
  \forall\,\{u,v\}\in E,\qquad u\in C\ \lor\ v\in C .
  \label{eq:cover}
\end{equation}
O objetivo da otimização é
\begin{equation}
  \OPT(G)=\min\bigl\{|C|:\ C\subseteq V
       \text{ satisfaz a Equação~\eqref{eq:cover}}\bigr\}.
  \label{eq:opt}
\end{equation}
Uma cobertura \emph{válida} satisfaz a condição de cobertura; uma cobertura
\emph{mínima} atinge o valor da Equação~\eqref{eq:opt}.
Já a \emph{melhor conhecida} é apenas a
menor cobertura válida encontrada durante uma busca. Ela só pode ser
declarada ótima quando há uma demonstração de que nenhuma cobertura menor
existe. Além disso, \emph{mínima} não significa apenas \emph{minimal}:
não conseguir remover individualmente um vértice de uma cobertura não
prova que sua cardinalidade é a menor possível.

\subsection{Classes de complexidade e NP-Completude}

As classes P e NP dizem respeito a problemas de decisão. P reúne os
problemas decididos em tempo polinomial por um algoritmo determinístico.
NP reúne os problemas cujas instâncias positivas possuem certificados de
tamanho polinomial verificáveis em tempo polinomial. Portanto, pertencer
a NP não significa, por si só, ser NP-completo \cite{garey1979,cormen2022}.

Uma redução polinomial $A\leq_p B$ transforma cada instância de $A$ em uma
instância de $B$, em tempo polinomial, preservando a resposta.
Um problema é NP-difícil quando todo problema de NP se reduz a ele.
Um problema de decisão é NP-completo quando, simultaneamente, pertence
a NP e é NP-difícil. A igualdade P=NP permanece em aberto; assim,
NP-Completude não autoriza afirmar que foi provada a inexistência de
algoritmos polinomiais \cite{garey1979}.

Para Vertex Cover de decisão, a entrada é $(G,k)$ e a pergunta é se existe
uma cobertura $C$ com $|C|\leq k$. Dado um conjunto candidato,
verificam-se seus índices, sua cardinalidade e as extremidades de todas
as arestas. Com marcações por vértice, essa verificação custa
$O(|V|+|E|)$, demonstrando o pertencimento a NP.

Para a NP-dificuldade, utiliza-se \emph{Independent Set}, problema de
decisão clássico NP-completo \cite{garey1979}. Dada a instância $(G,k)$,
com $0\leq k\leq |V|$, constrói-se a instância $(G,k')$ de Vertex Cover,
no mesmo grafo, em que
\begin{equation}
  k'=|V|-k,\qquad
  S\text{ é independente em }G
  \ \Longleftrightarrow\
  V\setminus S\text{ é cobertura de }G .
  \label{eq:reduction}
\end{equation}

Para provar o primeiro sentido da Equação~\eqref{eq:reduction},
suponha que $S$ seja independente.
Nenhuma aresta tem seus dois extremos em $S$; consequentemente, toda
aresta tem algum extremo em $V\setminus S$, que é uma cobertura.
Se $|S|\geq k$, essa cobertura tem tamanho no máximo $|V|-k$.
No sentido inverso, se $C$ é cobertura, nenhuma aresta pode ter ambos
os extremos em $V\setminus C$, pois ficaria descoberta. Logo,
$V\setminus C$ é independente e, se $|C|\leq |V|-k$, possui pelo menos
$k$ vértices.

A transformação preserva a resposta e é polinomial: mantém o grafo e
calcula um novo limite. Assim,
Independent Set $\leq_p$ Vertex Cover. Com o pertencimento a NP,
Vertex Cover de decisão é NP-completo. Essa equivalência se insere no
estudo clássico de reduções entre problemas combinatórios
\cite{karp1972,garey1979}. A versão de otimização é NP-difícil:
resolver o mínimo permitiria responder à versão de decisão por uma
comparação com $k$. Não se classifica aqui a otimização como
NP-completa.

\subsection{Aplicações e limites das modelagens}

Uma primeira modelagem possível é o \textbf{monitoramento de conexões em
redes de computadores}. Os vértices representam equipamentos, as arestas
representam conexões e os vértices selecionados recebem monitores.
Se observar uma das extremidades basta para monitorar a conexão,
uma cobertura representa uma seleção suficiente. Restrições de
capacidade, direção do tráfego, custo dos equipamentos e criptografia
não estão contempladas pelo modelo simples.

Uma segunda modelagem é a \textbf{seleção de pontos de inspeção em uma
malha de infraestrutura}. Vértices representam locais de inspeção e
arestas representam trechos entre esses locais. Se o requisito
operacional consiste em atribuir a cada trecho ao menos uma
extremidade com inspeção, minimizar os locais corresponde a Vertex
Cover. Isso não garante, por si só, que um posto inspecione toda a
extensão física de uma via, nem resolve restrições de deslocamento
ou de disponibilidade das equipes.

Uma terceira modelagem é a \textbf{instrumentação de uma rede de
componentes industriais}. Estações ou equipamentos são vértices e
interligações são arestas. Sob uma especificação que exija instrumentar
ao menos uma ponta de cada interligação, os vértices escolhidos indicam
onde instalar recursos. Se a aplicação real exigir ambas as pontas,
pesos diferentes, redundância ou cobertura espacial contínua, o problema
precisa ser reformulado. As três aplicações são correspondências de
modelagem, não estudos de caso validados pelos experimentos deste artigo.

\subsection{Abordagens algorítmicas existentes}

A enumeração por \textbf{força bruta} examina os $2^{|V|}$ subconjuntos,
verificando quais cobrem as arestas. É simples e pode servir como
oráculo em entradas pequenas, mas seu custo exponencial restringe
a escala. Técnicas de \textbf{ramificação} organizam essa enumeração
com escolhas locais; limites e regras de redução podem reduzir o
trabalho necessário, sem transformar a dificuldade geral em tempo
polinomial \cite{fomin2010}.

No \textbf{Branch and Bound}, uma solução factível limita superiormente
o mínimo procurado, enquanto limites inferiores permitem descartar
ramos sem potencial de melhoria. A implementação deste trabalho usa
ramificação pelos extremos de uma aresta e emparelhamento residual
para obter o limite inferior, conforme a Seção~\ref{sec:exact}.

A \textbf{aproximação por emparelhamento maximal} seleciona os extremos
das arestas de um matching maximal e garante fator~2
\cite{cormen2022,vazirani2001}. Outra estratégia possível é escolher
repetidamente o vértice de maior grau residual e remover suas arestas
incidentes. Essa \textbf{heurística gulosa por grau} privilegia a
cobertura imediata de muitas arestas. Não é o método aproximado
avaliado aqui, nem recebe garantia de fator~2 sem análise própria.

Também são relevantes formulações por \textbf{programação inteira},
com variáveis binárias $x_v$ e restrições $x_u+x_v\geq 1$ para cada
aresta. A relaxação linear admite arredondamento e métodos
primal-dual de aproximação \cite{vazirani2001}. Em grafos bipartidos,
a relação entre cobertura mínima e emparelhamento máximo permite
algoritmos polinomiais específicos \cite{cormen2022}; isso não
contradiz a NP-Completude em grafos gerais. O projeto experimental
implementa somente Branch and Bound e matching maximal, sem
resolvedores externos ou detecção especializada dessas famílias.
